\documentclass[11pt,a4paper]{article}
\usepackage[T1]{fontenc}
\usepackage[utf8]{inputenc}
\usepackage{lmodern}
\usepackage[margin=25mm]{geometry}
\usepackage{amsmath,amssymb,amsthm,mathtools,mathrsfs}
\usepackage{booktabs,longtable,array}
\usepackage{xcolor,microtype,enumitem,fancyhdr}
\usepackage{hyperref,url}
\definecolor{ink}{HTML}{16324F}
\definecolor{linkblue}{HTML}{145C90}
\hypersetup{colorlinks=true,linkcolor=linkblue,citecolor=linkblue,urlcolor=linkblue,
pdftitle={Direct quadratic block coding for Con(ZFC)},
pdfauthor={Research report prepared with ChatGPT for the ProveIt project}}
\setlength{\headheight}{14pt}
\pagestyle{fancy}
\fancyhf{}
\lhead{\small Direct quadratic block coding}
\rhead{\small Con(ZFC)}
\cfoot{\thepage}
\setlist{itemsep=3pt,topsep=5pt}
\newtheorem{theorem}{Theorem}[section]
\newtheorem{lemma}[theorem]{Lemma}
\newtheorem{proposition}[theorem]{Proposition}
\newtheorem{corollary}[theorem]{Corollary}
\theoremstyle{remark}
\newtheorem{remark}[theorem]{Remark}
\newcommand{\Z}{\mathbb Z}
\newcommand{\Qq}{\mathbb Q}
\newcommand{\N}{\mathbb N}
\newcommand{\Con}{\operatorname{Con}}
\newcommand{\B}{\mathcal B}
\newcommand{\Dc}{\mathcal D}
\newcommand{\Mc}{\mathcal M}
\newcommand{\Tc}{\mathcal T}
\newcommand{\Jc}{\mathcal J}
\newcommand{\Sc}{\mathcal S}
\newcommand{\Rc}{\mathcal R}
\renewcommand{\Cap}{\mathscr C}
\newcommand{\nocarry}{\mathrel{\perp_2}}
\newcommand{\file}[1]{\texttt{\path{#1}}}
\newcommand{\pin}{\texttt{22e75dd6b87b1b199bb08c778d7cf075ae5aae08}}
\allowdisplaybreaks[2]
\begin{document}
\hypersetup{pageanchor=false}
\begin{titlepage}
\vspace*{10mm}
{\color{ink}\Huge\bfseries Direct quadratic block coding\\[4mm]
for \(\Con(\mathrm{ZFC})\)\par}
\vspace{8mm}
{\Large A fixed equation in ten unrestricted integers\\
of exact degree \(237\,520\,456\)\par}
\vspace{10mm}
{\large Research report and reproducible construction\\
Prepared for Vladimir Reshetnikov and the ProveIt project\\
23 October 2026\par}
\vspace{10mm}
\noindent\begin{tabular}{@{}ll@{}}
\toprule
Unknowns & \(10\), all in \(\Z\)\\
Equation & One ordinary integer polynomial\\
Exact total degree & \(237\,520\,456\)\\
Binary arithmetic operations & \(809=222+25+562\)\\
Source & \(22\) quadratics in \(31\) positive integers\\
New coding positions & \(32\) pair-sum-distinct integers in \([0,799]\)\\
\bottomrule
\end{tabular}
\vspace{7mm}
\begin{abstract}
We replace the exponentially spaced scalar source coding in ProveIt's
ten-integer consistency equation by a direct encoding of its entire
quadratic residual system. A finite Sidon set distinguishes all quadratic
monomials. Reversed residual blocks share one carry test, while their
marked coefficients remain isolated even though the product blocks overlap.
This avoids both the quartic sum of squares and the powers-of-five
position scheme.

For the same fixed \(121\times8801\) compiler coefficient tuple as the
predecessor, the resulting polynomial has ten unrestricted integer
unknowns, exact degree \(237\,520\,456\), and an \(809\)-operation
division-free circuit. The predecessor has degree
\(589\,042\,901\,992\,797\,851\,562\,527\,064\) and \(1\,442\) operations.
The degree ratio is approximately \(2.48\times10^{18}\).

We prove the new source adapter, give the complete final polynomial by
finite arithmetic definitions, and compose it with the inherited compiler
endpoint and Sun's Lucas/interpolation theorems. The conclusion is
\(\mathrm{ZFC}\vdash [\,\forall\mathbf v\in\Z^{10},\
F_{10}(\mathbf v)\ne0\,]\leftrightarrow\Con(\mathrm{ZFC})\).
The archive includes coefficient bytes, deterministic generators,
independent algebraic checks, and exact-degree certificates. This is an
ordinary mathematical proof, not a new end-to-end proof-assistant
certificate or a claim of global minimality.
\end{abstract}
\vfill
{\small Audited against ProveIt commit\\\pin.}
\end{titlepage}
\hypersetup{pageanchor=true}
\tableofcontents
\newpage

\section{Result, comparison, and the meaning of simplicity}
\begin{theorem}[Main result]\label{thm:main}
The definitions in Sections \ref{sec:source} and \ref{sec:final},
with the fixed data in Appendix \ref{app:parent}, define a particular
polynomial
\[
F_{10}(f,g,h,k,l,w,x,y,z,m)\in
\Z[f,g,h,k,l,w,x,y,z,m]
\]
of exact total degree \(237\,520\,456\). Its delivered circuit has \(809\)
binary arithmetic operations. For the repository's specified proof
predicate, ordinary mathematics formalizable in ZFC proves
\begin{equation}\label{eq:main}
\exists(f,g,h,k,l,w,x,y,z,m)\in\Z^{10}\;F_{10}=0
\quad\Longleftrightarrow\quad \neg\Con(\mathrm{ZFC}).
\end{equation}
Consequently ZFC proves that \(F_{10}\) has no integer zero if and only if
\(\Con(\mathrm{ZFC})\).
\end{theorem}
All coefficients are fixed independently of the truth of consistency.
There is no unspecified recursively enumerable index, variable exponent,
radical, division, or hidden existential witness in the polynomial.
The intermediate names are polynomial abbreviations. They are not
additional unknowns.

\subsection{The comparison that is actually established}
The baseline is the structural-weight, merged-square, quintic
ten-integer equation in package 31 of the inspected collection
\cite{repo31}. We keep its whole coefficient tuple, its parent positive
system, and its final Lucas/norm construction. This isolates the effect
of the new source compiler.
\begin{table}[ht]
\centering\small
\begin{tabular}{@{}p{58mm}rr@{}}
\toprule
Measure & Package 31 & This construction\\
\midrule
Unrestricted integer unknowns & 10 & 10\\
Total degree & \(126496\cdot5^{31}+27064\) & \(237\,520\,456\)\\
Binary arithmetic operations & 1,442 & 809\\
Fixed circuit leaves (nonliteral) & 79 & 17\\
Weighted source records & 402 & 136\\
Exponent defining source-base constant & 4,269,882 & 2,134,940\\
Largest inherited compiler-port bit length & 1,070,004 & 1,070,004\\
\bottomrule
\end{tabular}
\caption{A comparison using the same four native compiler ports.
Constant-only recipes and reuse are free in both operation counts.}
\label{tab:comparison}
\end{table}
The exact old degree is
\[
589\,042\,901\,992\,797\,851\,562\,527\,064.
\]
Its ratio to the new degree is approximately \(2.4799670391\times10^{18}\).
The circuit saves \(633\) operations, a \(43.8974\%\) reduction.
The \(17\) fixed leaves consist of \(16\) coefficient recipes and one
explicit source-base constant.

These measures are not interchangeable. The collection also contains
positive-variable systems with far fewer arithmetic operations, including
the retained \(16\)-positive-variable, \(7\)-equation, degree-\(24\),
\(59\)-operation system, and the \(31\)-positive-variable quadratic
system with \(60\) operations \cite{repo}. Our equation does not improve
those operation counts or their maximum equation degrees. Its advance is
on the ten-\emph{unrestricted}-integer branch. We do not establish that
ten is minimal, that \(809\) is an optimal circuit size, or that the
displayed Sidon set is optimal.

\subsection{What is new and what is inherited}
The new contribution is the direct quadratic block adapter:
\begin{enumerate}
\item retain the \(22\) residual equations, without a sum of squares;
\item use pair-sum-distinct positions to distinguish their degree-two
monomials;
\item test all \(22\) selected coefficients in one binary no-carry
condition;
\item use block stride \(2M+1\), even though complete product blocks
overlap;
\item prove the size and positivity hypotheses for arbitrary bounded
\(g\), before decoding \(g\).
\end{enumerate}
Sidon sets and the Bose--Chowla construction are classical
\cite{bose,nathanson}. The binary-carry and binomial framework comes from
Sun \cite{sun2021}. The ten-integer Lucas/interpolation bridge comes from
Sun's 2026 preprint \cite{sun2026}. The quintic merger, coprime square
merger, separating weights, and compact norm formula are inherited from
package 31. The machine/compiler equivalence is inherited from ProveIt.
We prove a new adapter to those results, rather than assuming that Sun's
original positional lemma automatically permits different positions.

\section{Fixed source and dependency boundary}\label{sec:fixed}
\subsection{The actual source proposition}
Let \(R_1,\ldots,R_{22}\) be the quadratic residuals in the \(31\) positive
variables listed in Appendix \ref{app:parent}. The four fixed ports are
\[
(\beta_{\mathrm{native}},K_0,M_C,M_F)
=(\texttt{beta},\texttt{K},\texttt{MC},\texttt{MF\_source}).
\]
The high mask \(M_F\) is already shifted. All four ports come from the
same \(121\times8801\), \(943\)-tile compiler instance. Replacing only
some ports, or using the unshifted high mask, would specify a different,
unjustified instance.

The inherited positive-system/compiler theorem is
\begin{equation}\label{eq:inherited}
\begin{split}
\exists\mathbf v\in\Z_{>0}^{31}\ 
\bigwedge_{j=1}^{22}R_j(\mathbf v)=0
&\quad\Longleftrightarrow\quad
\text{K215 halts on its blank tape}\\
&\quad\Longleftrightarrow\quad
\neg\Con(\mathrm{ZFC}).
\end{split}
\end{equation}
The tuplewise \(32\)-to-\(31\) reduction, ordinary Pell/compiler proofs,
coefficient adapter, and endpoint reference are retained in
\file{vendor/}. In the upstream Lean source, the endpoint theorem is
\texttt{KEndpoint215.blank\_halts\_iff\_zfc\_refutation}; its conclusion
uses \texttt{BProv ZFCax\_s [] Form.fBot}. We have not rebuilt the full
Lean development in this delivery. Its inclusion records the exact
endpoint and does not turn the new polynomial reduction into a
kernel-checked theorem.

Shift every positive variable:
\[
\mathbf v=\mathbf n+\mathbf1,\qquad \mathbf n\in\N^{31},
\qquad \N=\{0,1,2,\ldots\}.
\]
Reorder the shifted residuals by putting original residual 5 last:
\[
\begin{split}
r_0,\ldots,r_{21}
={}&R_1(\mathbf n+\mathbf1),\ldots,R_4(\mathbf n+\mathbf1),
R_6(\mathbf n+\mathbf1),\ldots,\\
&R_{22}(\mathbf n+\mathbf1),R_5(\mathbf n+\mathbf1).
\end{split}
\]
This changes no common zero. Original residual 5 is \(E-X(a-E)\), whose
value at the all-one tuple is \(1\). Thus
\begin{equation}\label{eq:constant}
r_{21}(\mathbf0)=1,\qquad
\mathbf0\text{ is not a common zero}.
\end{equation}
The positive constant will also ensure positivity before digit decoding.

\subsection{Literal coefficients and a safe bound}
Write
\[
r_j(\mathbf n)=\sum_{|\alpha|\le2}a_{j,\alpha}\mathbf n^\alpha,\qquad
c_{j,\alpha}=(2-|\alpha|)!\prod_{i=1}^{31}\alpha_i!\;a_{j,\alpha}.
\]
Zero coefficients may be omitted. The generated file
\file{generated/source_records.json} lists all \(136\) nonzero
weighted records, including their exact finite recipes in the four
ports. An independent verifier reconstructs them from the pinned
parent sparse polynomial, with no call to the builder.
Set
\[
A_{\max}=\max_{j,\alpha}|c_{j,\alpha}|,\qquad
\beta_S=2^{2\,134\,940}.
\]
Exact integer evaluation of the delivered coefficient bytes gives
\begin{equation}\label{eq:bound}
\begin{gathered}
\operatorname{bitlength}(A_{\max})=2\,134\,926,\qquad
A_{\max}<2^{2\,134\,927},\\
\beta_S>8192 A_{\max}.
\end{gathered}
\end{equation}
The final exponent is even, so \(\beta_S\) is a power of \(4\).
This elementary, finite bound is checked as full integer arithmetic,
not inferred from modular samples. The large ports are compressed for
transport only: the polynomial's coefficients are the specified
uncompressed integers.

\section{A compact set of quadratic positions}\label{sec:sidon}
Use the following ordered set \(0=w_0<w_1<\cdots<w_{31}=799\):
\begin{equation}\label{eq:positions}
\begin{split}
(&0,10,35,55,62,105,144,145,159,168,232,245,306,323,325,353,\\
 &371,374,382,494,499,525,566,578,610,647,681,685,741,777,783,799).
\end{split}
\end{equation}
Every unordered pair sum \(w_i+w_k\), allowing \(i=k\), is distinct.
There are exactly \(\binom{33}{2}=528\) such sums. In fact they are
distinct modulo \(1023\), a stronger property than needed here.

The certificate is finite and transparent: enumerate all pairs and
check that the \(528\) resulting residues are distinct. Consequently
the main theorem does not depend on accepting an unimplemented
finite-field search or an optimality claim.

For provenance, the list was obtained from a Bose--Chowla construction
over \(\mathbf F_{2^{10}}\), represented by the polynomial
\(x^{10}+x^3+1\) and primitive element \(\theta=x\)
(encoded by the binary integer \(2\)), with the subfield of size
\(32\). A cyclic translate gives a diameter-\(901\) list. Multiplication
by \(47\), followed by translation by \(-342\), modulo \(1023\) gives
\eqref{eq:positions}. The independent Sidon verifier reconstructs that
provenance and checks the full primitive cycle. Only the final pair
certificate is needed for soundness.

The zero position represents the constant \(1\). It must be included:
checking only pair sums of the \(31\) positive positions would not
exclude collisions between constants, linear monomials, and quadratics.
For \(|\alpha|\le2\), write
\[
e(\alpha)=\sum_{i=1}^{31}\alpha_iw_i.
\]
Padding with \(2-|\alpha|\) copies of the zero position shows that
\(\alpha\mapsto e(\alpha)\) is injective on monomials of degree at most
two.

\section{The new simultaneous source adapter}\label{sec:source}
\subsection{Complete polynomial definitions}
Fix
\[
M=799,\quad s=2M+1=1599,\quad e_j=2M+js\quad(0\le j<22),
\]
and
\[
h_D=e_{21}=35177,\quad Q=h_D+2M=36775,\quad
t_N=M+Q+2=37576.
\]
Here \(h_D\) is a fixed integer, unrelated to the eventual unknown \(h\).
For the two integer unknowns \(f,g\), define, in the displayed order,
\begin{align}
b&=1+3f,&\B&=\beta_S b^2,&
\ell&=\sum_{i=1}^{31}\B^{w_i},\label{eq:base}\\
D_j(\B)&=\sum_{|\alpha|\le2}
c_{j,\alpha}\B^{2M-e(\alpha)},&
\Dc&=\sum_{j=0}^{21}\B^{js}D_j(\B),\label{eq:Dcode}\\
N_0&=\B^{M+1},&N_1&=4\B^{Q+1},&
N&=N_0N_1=4\B^{t_N},\label{eq:N}\\
\Mc&=N_0-1-(b-1)\ell,&
U_{\rm mask}&=(\B-2)\sum_{j=0}^{21}\B^{e_j},&
\Tc&=\Mc+N_0U_{\rm mask},\label{eq:mask}\\
\Jc&=2(1+g)^2\Dc+\sum_{i=0}^{Q}\B^{i+1},&
\Sc&=g+N_0\Jc,\label{eq:JS}\\
\Rc&=(\Sc+\Tc+1)N+\Tc+1,&
X&=(N-1)\Rc,&Y&=N^2,\label{eq:XY}\\
\Cap&=b\B^M.&&&&\label{eq:cap}
\end{align}
Each exponent is a fixed nonnegative integer. All sums are finite and
fixed. In particular these formulas are ordinary polynomials.
The expression \(U_{\rm mask}\) is distinct from the bridge polynomial
\(U\) in Section \ref{sec:final}.

\subsection{Why each target coefficient is exactly one residual}
\begin{lemma}[Simultaneous coefficient extraction]\label{lem:extract}
Let
\[
C_{\mathbf n}(T)=\left(1+\sum_{i=1}^{31}n_iT^{w_i}\right)^2,\qquad
D(T)=\sum_{j=0}^{21}T^{js}
\sum_{|\alpha|\le2}c_{j,\alpha}T^{2M-e(\alpha)}.
\]
For every \(j\), the coefficient of \(T^{e_j}\) in
\(C_{\mathbf n}(T)D(T)\) is \(2r_j(\mathbf n)\).
\end{lemma}
\begin{proof}
The coefficient of \(T^{e(\alpha)}\) in \(C_{\mathbf n}\) is
\[
\frac{2!}{(2-|\alpha|)!\prod_i\alpha_i!}\,\mathbf n^\alpha.
\]
This uses the pair-sum property including the zero position. A term
from block \(k\) can contribute to target \(e_j\) only if
\[
e(\gamma)-e(\alpha)=(j-k)s.
\]
The left side lies in \([-2M,2M]\), while \(s=2M+1\).
Therefore \(k=j\). Within that block, injectivity gives
\(\gamma=\alpha\). Multiplying the two coefficients cancels the
factorials and gives \(2a_{j,\alpha}\mathbf n^\alpha\).
Summing proves the claim.
\end{proof}
This is where two separate savings occur. Quadratic monomials require
only a \(B_2\) set, whereas squaring the residuals first would require
degree-four separation. Further, full product blocks need not be
disjoint. The following intervals explain the shorter stride:
\[
\begin{array}{c|c|c}
\text{block }j & \text{reversed source support} & \text{product support}\\
\hline
j &[js,js+2M]&[js,js+4M]
\end{array}
\]
The marked target is \(js+2M\). The preceding product block ends at
that target minus \(1\); the next product block starts at that target
plus \(1\). Overlap elsewhere is harmless.

\subsection{Bounds before decoding: avoiding a circular proof}
\begin{lemma}[Global source bounds]\label{lem:global}
Suppose \(f\ne0\), \(b=1+3f\) is an integer square, and
\(0\le g<2\Cap\). Then \(f>0\), \(b\ge4\), and
\[
0\le g<N_0,\quad \Dc>0,\quad
0<\Jc<2\B^{Q+1}<N_1,\quad 0\le\Sc,\Tc<N.
\]
Moreover \(3\mid X\), \(X>3b\), and \(Y>\max\{b,256\}\).
These statements do not assume that \(b\) is a power of two or that
\(g\) has the intended digit representation.
\end{lemma}
\begin{proof}
A square congruent to \(1\) modulo \(3\) is positive; excluding \(f=0\)
excludes \(b=1\), so \(b\ge4\). Also \(\B>2b\), hence
\(g<2b\B^M<\B^{M+1}=N_0\).

The reversed source blocks in \(\Dc\) occupy disjoint intervals of
length \(2M+1\). Each coefficient has absolute value at most \(A_{\max}\).
By \eqref{eq:constant}, the leading coefficient, at exponent \(h_D\),
is exactly \(2\). Since \(\B-1>A_{\max}\),
\[
\Dc\ge 2\B^{h_D}
-A_{\max}\frac{\B^{h_D}-1}{\B-1}>\B^{h_D}>0.
\]
The upper bound
\[
\Dc\le A_{\max}\sum_{i=0}^{h_D}\B^i
<\frac{A_{\max}\B^{h_D+1}}{\B-1}
\]
also holds. Since \(g\) is integral, \(1+g\le2b\B^M\).
Consequently
\[
0<\Jc<
\frac{8A_{\max}b^2+\B}{\B-1}\B^{Q+1}
<2\B^{Q+1}.
\]
The final inequality follows from \(\B=\beta_Sb^2\) and
\(\beta_S>8192A_{\max}\).

The base-\(\B\) digits of \(\Mc\) are \(\B-b\) at source
positions and \(\B-1\) at all other positions through \(M\).
Thus \(0\le\Mc<N_0\). Also
\[
0\le U_{\rm mask}<\B^{Q+1}<N_1.
\]
The asserted bounds on \(\Sc,\Tc\) follow by concatenation.
In particular \(\Rc\ge N+1>\B\ge b\); the fixed large \(\beta_S\)
makes \(X>3b\) and \(Y>256\) immediate.

Finally \(b\equiv1\pmod3\), \(\beta_S\equiv1\pmod3\), and
\(N=4\B^{t_N}\equiv1\pmod3\). Thus \(3\mid N-1\), so \(3\mid X\).
\end{proof}
The positive final block is useful for a specific reason: it makes
\(\Dc>0\) for every bounded candidate \(g\). Positivity inferred only
after decoding would not justify the hypotheses used to obtain that
decoding.

\subsection{Low digits and all residuals in one carry condition}
For nonnegative integers \(u,v\), write \(u\nocarry v\) when their
binary addition has no carry; equivalently, their sets of nonzero
binary positions are disjoint.
\begin{lemma}[Decoding the combined mask]\label{lem:mask}
Assume the hypotheses of Lemma \ref{lem:global}, and assume additionally
that \(b\) is a power of two. Then
\[
\Sc\nocarry\Tc
\quad\Longleftrightarrow\quad
\left\{
\begin{array}{l}
g=\displaystyle\sum_{i=1}^{31}n_i\B^{w_i}
\text{ for some }0\le n_i<b,\\
r_j(\mathbf n)=0\quad(0\le j<22).
\end{array}\right.
\]
\end{lemma}
\begin{proof}
Now \(\B,N_0,N_1,N\) are powers of two. The global bounds make the
low and high parts genuine, nonoverlapping binary blocks, so
\[
\Sc\nocarry\Tc
\iff (g\nocarry\Mc)\ \text{and}\ (\Jc\nocarry U_{\rm mask}).
\]
At a source position the low mask digit is \(\B-b\), with all the
binary positions above the lowest \(\log_2 b\) positions occupied.
At every other base position its digit is \(\B-1\). Thus
\(g\nocarry\Mc\) gives exactly the displayed representation with
\(0\le n_i<b\). In particular the constant position zero is forbidden.

On this decoded branch, let \(\theta_i\) be the coefficients of
\(C_{\mathbf n}(T)D(T)\). The sum of the nonnegative coefficients
of \(C_{\mathbf n}\) is
\[
\left(1+\sum_i n_i\right)^2<(31b)^2.
\]
Because the coefficients of \(D\) have absolute value at most
\(A_{\max}\), convolution gives
\[
|\theta_i|<961A_{\max}b^2<\B/2.
\]
Therefore
\[
\Jc/2=\sum_{i=0}^{Q}(\B/2+\theta_i)\B^i
\]
has genuine base-\(\B\) digits strictly between \(0\) and \(\B\).
Division by two occurs only in this proof; the defining polynomial
\(\Jc\) is integral and division-free.

Doubling both summands preserves the absence of binary carries.
Thus the high condition is equivalent to
\[
\Jc/2\nocarry
(\B/2-1)\sum_{j=0}^{21}\B^{e_j}.
\]
At each marked position the mask occupies every binary position
below the top bit. A digit in \(\{1,\ldots,\B-1\}\) avoids this
mask if and only if it is exactly \(\B/2\). Hence
\(\theta_{e_j}=0\) for every \(j\), and Lemma \ref{lem:extract}
identifies this as \(2r_j(\mathbf n)=0\).
\end{proof}

\subsection{The binomial interface}
\begin{lemma}[Binary binomial test]\label{lem:binomial}
If \(N=2^u\) and \(0\le S,T<N\), put
\[
X=(N-1)((S+T+1)N+T+1),\qquad Y=N^2.
\]
Then
\[
Y\mid\binom{2X}{X}\quad\Longleftrightarrow\quad S\nocarry T.
\]
No hypothesis \(S+T<N\) is necessary.
\end{lemma}
\begin{proof}
Write \(s_2(v)\) for the number of nonzero binary digits of \(v\).
The exact identity
\[
X=(S+T)N^2+(N-1-S)N+(N-1-T)
\]
separates the lower two blocks, even when \(S+T\ge N\).
Hence
\[
s_2(X)=2u+s_2(S+T)-s_2(S)-s_2(T).
\]
The factorial valuation identity gives
\(v_2\binom{2X}{X}=s_2(X)\).
Binary addition gives \(s_2(S+T)\le s_2(S)+s_2(T)\), with
equality exactly when there is no carry. Divisibility by \(2^{2u}\)
is therefore equivalent to that equality.
\end{proof}
This is the binary form of the carry mechanism used by Sun
\cite[Lemmas 2.2 and 2.4]{sun2021}; the displayed proof specifies
the exact normalization needed here.

\begin{proposition}[Source interface]\label{prop:interface}
The polynomials \eqref{eq:base}--\eqref{eq:cap} have the following
properties.
\begin{enumerate}
\item If the \(r_j\) have a common zero in \(\N^{31}\), then for
arbitrarily large positive integers \(f\) there is an integer \(g\)
with
\[
b\le g<\Cap,\qquad b\text{ a square and a power of }2,\qquad
Y\mid\binom{2X}{X}.
\]
\item Conversely, if \(f\ne0\), \(0\le g<2\Cap\), \(b\) is a square
and a power of \(2\), and \(Y\mid\binom{2X}{X}\), then the \(r_j\)
have a common zero in \(\N^{31}\).
\item Whenever \(f\ne0\), \(b\) is a square, and
\(0\le g<2\Cap\), the size bounds and divisibility \(3\mid X\)
in Lemma \ref{lem:global} hold. The cap has the fixed form
\[
\Cap=2^{2\,134\,940M}b^{2M+1}.
\]
\end{enumerate}
\end{proposition}
\begin{proof}
For completeness, choose arbitrarily large \(b=4^a\) greater than
all \(n_i\), set \(f=(b-1)/3\), and encode
\(g=\sum_i n_i\B^{w_i}\).
Equation \eqref{eq:constant} excludes the zero tuple. All allowed
positions are positive, so \(g\ge\B\ge b\). Also
\[
g\le(b-1)\sum_{i=1}^{31}\B^{w_i}
\le(b-1)\sum_{i=0}^{M}\B^i<b\B^M=\Cap;
\]
the last inequality follows from \(\B>b\).
Apply Lemmas \ref{lem:mask} and \ref{lem:binomial}.
The reverse implication is those same lemmas in reverse order,
using the global bounds first. The last assertion follows directly
from Lemma \ref{lem:global} and the definition of \(\B\).
\end{proof}

\section{The complete ten-integer equation}\label{sec:final}
We now retain the package-31 finalizer. This section gives every
formula so that the equation can be read without executing code.
The ten unrestricted integer unknowns are exactly
\[
\boxed{\,f,g,h,k,l,w,x,y,z,m\,}.
\]
Neither the \(31\) source witnesses nor any of the abbreviations below
are quantified in the final equation.

\subsection{Lucas bridge}
Use \(b,X,Y,\Cap\) from Section \ref{sec:source}, then define
\begin{align}
L&=lY,&U&=2LX,&V&=4gwY,&W&=bw,\nonumber\\
A&=U(V+1),&B&=2X+1,&C&=B+(A-2)h,\nonumber\\
P&=U^2V,&K&=X+1+k(P-2),\label{eq:bridge}\\
S&=2A-5,&T&=3WC-2(W^2-1),\nonumber\\
R_2&=(P^2-4)K^2+4,\nonumber
\end{align}
\[
\Delta=8\Cap^3gK^2-g^2\{32(C-KL)^2\Cap^3+g^2K^2\},
\qquad O=f^2l^2\Delta.
\]
Here the ordinary letter \(B\) is different from the radix \(\B\).
The unknown \(w\) is different from the fixed positions \(w_i\).
The odd integer \(S=2A-5\) is always nonzero.
Next define
\begin{align}
\xi&=zx,&D&=(A^2-4)C^2+4,&E&=C^2D\xi,\nonumber\\
F&=4(A^2-4)E^2+1,&
G^*&=2(A+2)(A-2)^2E^2-1,\label{eq:lucas}\\
\Omega&=16E,&H^*&=C+(B-C)F+\Omega Fy,\nonumber\\
I^*&=((G^*)^2-1)(H^*)^2+1.\nonumber
\end{align}
These are the formulas of Sun's modified Lucas certificate and
binomial bridge \cite[Theorems 3.1 and 4.1; Propositions 6.1 and
6.2]{sun2026}.

For clarity, the precise imported interface is as follows.
The forward bridge, applied to Proposition \ref{prop:interface},
produces \(f,g,h,k,l,w\) with \(b\) square, \(R_2\) square,
\(S\mid T\), \(C=u_B(A)\), and \(O>0\), where
\(u_0(A)=0,u_1(A)=1,u_{n+1}(A)=Au_n(A)-u_{n-1}(A)\).
One may choose \(f\) arbitrarily large. On this branch \(g\ge b\ge4\);
\(\Delta\) is a positive multiple of \(g\), so \(O\ge4\).
The interpolation theorem permits a positive \(\xi\) divisible by a
prescribed positive integer and a \(y\) in any prescribed residue
class modulo \(\xi\), while making \(DFI^*\) square.

In the reverse direction, \(b\) square, \(O>0\), \(\xi\ne0\),
\(R_2\) square, \(DFI^*\) square, and \(S\mid T\) imply the source
conditions of Proposition \ref{prop:interface}(2).
The order of the bounds is explicit: \(O>0\) first gives
\(f,l\ne0\) and \(\Delta>0\). The square \(b\) then gives
\(b\ge4\) and \(\Cap>0\). The displayed formula for \(\Delta\)
forces \(K\ne0\) and \(g>0\); dropping its other nonnegative
subtracted term gives \(g^3<8\Cap^3\), hence \(g<2\Cap\).
Lemma \ref{lem:global} is therefore available before decoding.
The Lucas bridge then recovers the required power-of-two and
binomial conditions.
Inspection of the cited proofs shows that their uses of the source
are exactly its bounds, cap shape, forward and reverse implications,
and \(3\mid X\). They do not require a recurrence among the coding
positions. This is the reason the new adapter composes with them.
We invoke these number-theoretic results, rather than claiming to
reprove their Lucas-sequence arguments.

\subsection{Quintic divisibility and nonvanishing merger}
Set
\begin{align}
\pi(v)&=(v^3-v-1)(v^2+23),&\Sigma&=\xi S^5,\nonumber\\
r&=T+Sy,&
\Theta&=(r^3-rS^2-S^3)(r^2+23S^2).\label{eq:quintic}
\end{align}
\begin{lemma}\label{lem:quintic}
The polynomial \(\pi\) has no rational zero and has a zero modulo
every positive integer. If \(S\ne0\) and \(\xi\ne0\), then
\(\Sigma\mid\Theta\) implies \(S\mid T\).
If \(S\mid T\), some residue class of \(y\) modulo \(\xi\)
makes \(\Sigma\mid\Theta\).
Moreover \(\Theta\ne0\) for all integer \(T,y\) and nonzero \(S\).
\end{lemma}
\begin{proof}
The cubic is monic and its only possible rational roots are
\(\pm1\), neither a root. The quadratic is positive on \(\Qq\).
For an odd prime \(p\ne23\), the cubic has discriminant \(-23\).
If it has a root modulo \(p\), that root is simple. Otherwise it is
irreducible of degree three, whose Frobenius permutation is an even
3-cycle; its discriminant is then a square modulo \(p\). The quadratic
has a simple root. At \(p=23\), the cubic has the simple root \(3\).
Hensel lifting gives roots modulo odd prime powers.
Since \(-23\equiv1\pmod8\), it is a square in \(\Z_2\), giving
roots modulo all powers of \(2\). The Chinese remainder theorem
gives the asserted local roots.

If \(S\nmid T\), choose a prime \(p\) with
\(\alpha=v_p(S)>\beta=v_p(T)\).
Then \(v_p(r)=\beta\). The first homogeneous factor has valuation
\(3\beta\) and the second \(2\beta\), because their other terms
have strictly larger valuation. Thus
\(v_p(\Theta)=5\beta<5\alpha\), contradicting \(S^5\mid\Theta\).

Conversely write \(T=Sq\), choose a root \(c\) of \(\pi\) modulo
\(|\xi|\), and prescribe \(y\equiv c-q\pmod{\xi}\).
Then \(\Theta=S^5\pi(q+y)\) is divisible by \(\xi S^5\).
Finally, \(\Theta=0\) would give the rational root \(r/S\) of
\(\pi\), which is impossible.
\end{proof}

\subsection{One square condition for positivity and the square base}
Define
\begin{equation}\label{eq:radicands}
\rho_0=b\{1+(3O-4)b^2z^2\},\qquad
\rho_1=R_2,\qquad \rho_2=DFI^*.
\end{equation}
Suppose \(z\ne0\) and \(\rho_0\) is an integer square.
The factors \(b\) and \(1+(3O-4)b^2z^2\) are coprime.
A nonzero product that is a square forces both to be signed squares
of the same sign. But \(b\equiv1\pmod3\) cannot be a negative square.
The zero-product case likewise forces \(|b|=1\), hence \(b=1\).
Therefore \(b\) is a positive square and the bracket is a
nonnegative square. Since \(bz\ne0\), an integer \(O\le0\)
would make that bracket at most \(-3\). Thus \(O>0\).

Conversely, on the positive bridge branch \(O\ge4\) and \(b\) is a
square. The positive integer \(d=3O-4\) is \(2\) modulo \(3\), so
it is not a square. A positive Pell unit for
\(u^2-dv^2=1\) can be raised to an exponent making its residue
equal to \(1\) modulo \(b\). This yields a solution with \(v>0\)
divisible by \(b\). Write \(v=bz\). Then \(z>0\) and \(\rho_0\)
is square. Sun's interpolation theorem may now take this \(z\)
as its prescribed divisibility parameter for \(\xi\); set
\(x=\xi/z\). Thus the same existing unknown \(z\) serves both
purposes.

\subsection{Separating weights and a radical-free norm}
Whenever \(z\ne0\), the actual polynomial \(\rho_0\) satisfies
\(|\rho_0|>1\). If \(\rho_0=0\), then
\((3O-4)b^2z^2=-1\), forcing \(|b|=1\), hence \(b=1\), \(f=0\),
and \(O=0\), a contradiction. If \(\rho_0=\pm1\), the same
argument gives \(\rho_0=1-4z^2\le-3\).
Also
\[
\sqrt{|\rho_1|}\le P^2+K^2+2,
\]
because \(|\rho_1|\le(P^2+4)K^2+4\) and the difference from
\((P^2+K^2+2)^2\) is \(P^4+K^4+P^2K^2+4P^2\ge0\).
Set
\begin{equation}\label{eq:weights}
v_0=1,\qquad v_1=\rho_0,\qquad v_2=\rho_0(P^2+K^2+3).
\end{equation}
They satisfy
\[
|v_1|>\sqrt{|\rho_0|},\qquad
|v_2|>\sqrt{|\rho_0|}+|v_1|\sqrt{|\rho_1|}.
\]
In the multiquadratic field generated by the square roots of the
\(\rho_i\), a rational signed sum
\(\sum_i\varepsilon_i v_i\sqrt{\rho_i}\) forces every \(\rho_i\)
to be an integer square. Indeed, if an automorphism changes a
nonzero root term, take the largest changed index. Its magnitude
is at least \(|v_i|\), larger than the sum of all earlier
magnitudes; its change cannot cancel. Hence each root is rational,
and a rational square root of an integer is integral. Negative
radicands are covered by complex absolute values; zero radicands
already are squares.

Finally define
\begin{align}
Z&=\Sigma m-\Theta,&
A_i&=\Sigma^2v_i^2\rho_i\quad(0\le i\le2),\nonumber\\
V_n&=Z^2+A_2-A_0-A_1,&
W_n&=V_n^2+4Z^2A_2-4A_0A_1.\label{eq:normparts}
\end{align}
The promised equation is
\begin{equation}\label{eq:final}
\boxed{\qquad F_{10}=W_n^2-16V_n^2Z^2A_2=0.\qquad}
\end{equation}
Everything in \eqref{eq:final} is defined by ordinary polynomial
arithmetic above. The proof identity is
\[
F_{10}=\prod_{\varepsilon\in\{\pm1\}^3}
\left(Z+\sum_{i=0}^2\varepsilon_i\sqrt{A_i}\right).
\]
The radicals occur only in this identity, not in the definition of
the equation. Expanding the formal norm gives \(35\) terms.

\subsection{Both directions}
If the parent positive system has a solution, Proposition
\ref{prop:interface} and the forward bridge give
\(f,g,h,k,l,w\), with \(O\ge4\), \(b,\rho_1\) squares,
\(S\mid T\), and the Lucas identity. The Pell construction gives
a positive \(z\) making \(\rho_0\) square. Interpolation supplies
\(\xi>0\) divisible by \(z\), and a \(y\) in the residue class
specified by Lemma \ref{lem:quintic}, with \(\rho_2\) square.
Set \(x=\xi/z\). Then \(\Sigma\ne0\) and \(\Sigma\mid\Theta\).
Choosing
\[
m=\Theta/\Sigma-\sum_i\varepsilon_i v_i\sqrt{\rho_i}
\]
gives an integer \(m\) making one norm factor zero.

Conversely suppose \(F_{10}=0\). Since \(S\ne0\), Lemma
\ref{lem:quintic} gives \(\Theta\ne0\).
If \(\Sigma=0\), all \(A_i=0\) and \eqref{eq:final} becomes
\(\Theta^8\ne0\). Therefore \(\Sigma\ne0\), so
\(\xi=zx\ne0\) and \(z\ne0\). A vanishing norm factor, divided
by \(\Sigma\), gives a rational weighted sum of the roots
of \(\rho_i\). The separating-weight argument makes all three
radicands integer squares. That sum is now integral, so
\(\Sigma\mid\Theta\). Hence \(S\mid T\), while the square
\(\rho_0\) gives \(b\) square and \(O>0\).
The reverse bridge and Proposition \ref{prop:interface} yield a
common nonnegative zero of the shifted residuals.
Composing with \eqref{eq:inherited} proves Theorem \ref{thm:main}.

This proof does not assume that ZFC is consistent, and does not
decide whether the displayed equation has a solution. In particular
one must not replace the equivalence by an unconditional assertion
of no solutions. If ZFC is consistent and the standard hypotheses
of the second incompleteness theorem hold, ZFC cannot prove
this equation has no integer zero. Consistency alone should not
be advertised here as a proof of both directions of independence.

\section{Exact degree and circuit size}\label{sec:degree}
\subsection{A useful degree formula}
The source proof applies to \(m_r\) quadratic residuals in the same
\(31\) nonnegative coordinates, on a pair-sum-distinct set of
maximum \(M\), with the same normalization
and coefficient bounds. Then
\[
Q=(m_r-1)(2M+1)+4M,\qquad
t_N=M+Q+2=(2m_r+3)M+m_r+1.
\]
For the present bridge, direct degree propagation gives
\[
\deg F_{10}\le6320t_N+48M+1784.
\]
For \(m_r=22\), this is
\[
297088M+147144.
\]
At \(M=799\), it equals \(237\,520\,456\).
This explains quantitatively why reducing the source positions
matters so much.

For this specific coefficient tuple the relevant leading terms
are nonzero. The following table records the main degree
propagation, with \(t=t_N\):
\begin{longtable}{@{}lll@{}}
\toprule
Polynomial & Degree expression & Value\\
\midrule
\endhead
\(N\) & \(2t\) & 75,152\\
\(X\) & \(6t\) & 225,456\\
\(Y\) & \(4t\) & 150,304\\
\(\Cap\) & \(2M+1\) & 1,599\\
\(A\) & \(14t+3\) & 526,067\\
\(C\) & \(14t+4\) & 526,068\\
\(P\) & \(24t+4\) & 901,828\\
\(K\) & \(24t+5\) & 901,829\\
\(R_2\) & \(96t+18\) & 3,607,314\\
\(\Delta\) & \(56t+6M+17\) & 2,109,067\\
\(O\) & \(56t+6M+21\) & 2,109,071\\
\(\rho_0\) & \(56t+6M+26\) & 2,109,076\\
\(D\) & \(56t+14\) & 2,104,270\\
\(E\) & \(84t+24\) & 3,156,408\\
\(F\) & \(196t+54\) & 7,364,950\\
\(G^*\) & \(210t+57\) & 7,891,017\\
\(H^*\) & \(280t+79\) & 10,521,359\\
\(I^*\) & \(980t+272\) & 36,824,752\\
\(\rho_2\) & \(1232t+340\) & 46,293,972\\
\(v_2\) & \(104t+6M+36\) & 3,912,734\\
\(\Sigma\) & \(70t+17\) & 2,630,337\\
\(A_2\) & \(1580t+12M+446\) & 59,380,114\\
\(F_{10}\) & \(6320t+48M+1784\) & 237,520,456\\
\bottomrule
\end{longtable}
In particular \(A_2\) has strictly greater degree than \(A_0,A_1\)
and \(Z^2\). The final leading term is \(A_2^4\), with no
competitor of its degree. The audit also certifies degree
attainment independently by substituting a single indeterminate
for all ten unknowns and tracking its leading coefficient modulo
three primes:
\[
\begin{array}{c|rrr}
\text{prime}&101&1009&1000003\\\hline
\text{coefficient of degree }237520456&80&105&684130
\end{array}
\]
A single nonzero entry suffices to prove that the circuit upper
bound is attained. This does not require expansion of an
astronomically large multivariate polynomial.

\subsection{What the operation count charges}
The delivered directed acyclic circuit has
\[
222\text{ binary additions},\quad
25\text{ binary subtractions},\quad
562\text{ binary multiplications}.
\]
Every multiplication by a literal integer is charged unless
removed by a constant identity. Fixed powers are expanded by
binary powering; fixed finite geometric sums use
\[
G_{2a}(B)=G_a(B)(1+B^a),\qquad
G_{2a+1}(B)=G_{2a}(B)+B^{2a}.
\]
No division node occurs. Reuse of a previous node is free.
Constant-only arithmetic is folded, as in the predecessor.
The \(16\) fixed recipes are finite integer formulas in the four
provided ports; the seventeenth leaf is \(2^{2134940}\).

These are arithmetic-circuit counts, not bit-operation counts.
They do not include generating the million-bit coefficients,
and they say nothing about the size of a hypothetical solution.
The circuit is an exact finite specification of an ordinary
polynomial, even though printing its full expanded coefficient
list would be an unsuitable artifact.

\section{Verification, reproducibility, and proof status}\label{sec:checks}
\subsection{Independent checks supplied}
The main verifier reads the delivered circuit, sparse parent,
coefficient recipes, and compressed port bytes. It does not import
the circuit builder. The retained receipts report:
\begin{itemize}
\item all \(528\) pair sums distinct, including the zero position;
\item all \(22\) shifted quadratics reconstructed symbolically;
\item all \(136\) factorial-weighted coefficient records verified;
\item all \(22\) target-coefficient identities verified as exact
polynomial identities in the source variables and fixed ports,
including every possible cross-block contribution;
\item the four active raw port hashes and coefficient size bound
verified by full integer arithmetic;
\item \(76\) named circuit nodes compared against an independent
formula transcription at \(48\) assignments over three finite fields;
\item exact degree attained, with the three nonzero leading
coefficients displayed above;
\item complete small-system source checks, invalid low-digit
rejections, a deliberate colliding-position negative control,
and arbitrary-\(g\) upper-boundary cases;
\item regeneration of the native coefficient tuple from the literal
transition and tile data, with comparison to the active port hashes.
\end{itemize}
The symbolic target checks are exhaustive identities. The finite-field
node comparisons are transcription tests, not proofs of the unbounded
Lucas or compiler statements. The small examples likewise test the
source implementation without deciding the existence of an integer
zero of the final equation. The main source equivalence is established
by the proof in Section \ref{sec:source}.

\subsection{Reproducing the archive}
The top-level \file{README.md} is the operational entry point.
Run
\begin{quote}\ttfamily
python3 reproduce.py
\end{quote}
to regenerate the circuit and its independent receipts with the
Python standard library. Run
\begin{quote}\ttfamily
latexmk -pdf -interaction=nonstopmode -halt-on-error report.tex
\end{quote}
to build the article. The compiler-port reproduction instructions
are in \file{vendor/README.md}. All active data and build inputs
are local to the package; no network request or solver query is
needed for arithmetic reproduction.

\file{INSTANCE.json} identifies the authoritative binding and
polynomial. \file{generated/signed10_b2.json} is the complete
live-node circuit. \file{generated/source_records.json} contains
the coefficients needed to interpret every fixed leaf.
\file{MANIFEST.sha256} records the delivered files.
Proof notes retain the independent source reviews, while
\file{vendor/} records the inherited material and attribution.

The repository baseline is pinned at \pin.
Some inherited files name earlier commits for their own sources;
these are provenance records, not instructions to mix coefficient
tuples. In particular any default binding metadata in the raw
parent sparse file is superseded by the explicit whole-tuple
binding in this package.

\subsection{What remains outside the certificate}
The new source adapter and its algebraic instantiation are given
in full. The semantic machine/compiler chain and Sun's 2026
Lucas/interpolation theorems remain named mathematical
dependencies. The latter source is a 2026 preprint; this report
does not supply a new formal verification of it.
The archive includes the endpoint reference and inherited proof
material, but not a rebuild of the entire ProveIt proof-assistant
project. Thus the correct status is an ordinary proof with
explicit dependencies and extensive finite algebraic verification.

No conclusion here establishes the truth or falsity of
\(\Con(\mathrm{ZFC})\). Nor do finite tests establish a new
independence theorem. They certify the explicitly stated arithmetic
construction and its transcription.

\section{Further reductions suggested by the construction}
The practical change is to optimize the source compiler before
optimizing the final norm. It suggests several concrete directions.

\paragraph{Support-aware positions.}
A full \(B_2\) set separates every quadratic monomial, whereas the
actual \(22\) quadratics use only a small subset. It may suffice to
exclude collisions between a source-supported monomial and any
multiplier monomial that could reach its target. A smaller list
could lower \(M\), and the present degree formula would lower the
degree by \(297088\) for each unit saved. Such a construction needs
an explicit collision certificate; simply dropping unused pairs
without auditing all multiplier monomials is unsound.

\paragraph{Support-aware targets.}
Uniform stride \(2M+1\) is a simple, proved sufficient condition.
The individual reversed blocks often have sparse or shorter
support. Nonuniform target positions could reduce their total
height, provided every cross-block contribution to a marked
coefficient is excluded. The current independent symbolic
extraction checker is a suitable certificate format for that
search.

\paragraph{Better small source systems.}
Reducing the number of source variables changes the size of the
required additive set; reducing the number of residuals shortens
the block range. These objectives may trade against source
degree, coefficient size, and evaluation cost. A smaller
positive system cannot be substituted without its complete
sign, normalization, and compiler-interface proof.

\paragraph{Higher-degree simultaneous coding.}
The coefficient-extraction argument generalizes to a \(B_d\) set,
with the zero position representing padding to degree \(d\).
Weights \((d-|\alpha|)!\prod_i\alpha_i!\) extract
\(d!\,r_j\); stride \(dM+1\) isolates targets.
This is an algebraic generalization of Lemma \ref{lem:extract}.
A complete higher-degree source theorem must separately redo
the coefficient, carry, and reverse-branch bounds; it should not
be asserted solely from the algebraic identity.

\paragraph{Formalization.}
The cleanest next formal target is the source interface:
finite monomial separation, the target-coefficient identity,
binary masks, and the binomial valuation lemma. These are
independent of K215 and can be stated as reusable arithmetic
lemmas. The existing circuit and record format then permits a
small certificate checker. Formalizing the complete result
also requires the inherited Lucas and compiler dependencies.

\paragraph{Nine variables and global optimization.}
The result does not give a nine-unrestricted-integer equation.
Nor does it automatically improve the repository's
nine-positive-variable construction: that uses a different
compressor whose interface must be re-established.
Such extensions are research tasks, not consequences of a
variable-name substitution.

\clearpage
\appendix
\section{The pinned positive source system}\label{app:parent}
The \(31\) strictly positive variables, in the order used by the
coefficient records and weights \(w_1,\ldots,w_{31}\), are
\[
\begin{split}
(&J,F,\alpha,z,f,h,i,o,j,s,w,\tau,\eta,\zeta,y,\gamma,\\
 &N,E,P,c,a,X,D,C_2,L,F_2,V,g,I,B_{\rm aux},\Delta).
\end{split}
\]
These names belong to the \emph{parent} system and are local to
this appendix; they are not identified with the final ten
unknowns. Define the affine abbreviations
\[
q=\beta_{\rm native}J+1,\qquad k=\eta+\zeta,\qquad
H=4a+3,\qquad C=q-F-\alpha.
\]
The \(22\) equations are
\begin{align*}
L&=q^2-1,\\
N&=q(L+1),\\
X&=wq,\\
a&=E+sN,\\
E&=X(a-E),\\
c&=k(a-E)+\eta,\\
P&=E(c-\eta)+k,\\
\tau^2&=P(P-k)+1,\\
D&=X+ac+\gamma H,\\
\Delta&=a^2+H,\\
C_2&=c^2,\\
D^2&=\Delta C_2+1,\\
k&=jc-V+1+hE,\\
jc-V&=(q-F)(N-q)-CL+(M_C+qM_F)J,\\
(K_0+w)C+(q-F)-z(q-1)&=1,\\
F_2&=f^2,\\
V&=of-c,\\
g&=y^2-V^2,\\
I&=iC_2,\\
B_{\rm aux}&=I^2,\\
B_{\rm aux}&=\Delta(F_2-1),\\
y^2-B_{\rm aux}g&=1.
\end{align*}
Residuals mean left side minus right side after the affine
abbreviations are expanded. They are quadratic in these \(31\)
variables, with the four compiler ports treated as constants.
The copied sparse file has SHA-256
\begin{quote}\small\ttfamily
1c336ba0b3c9827ffddcf7ce895cad9dca2907f907cde7484c7f387bb2cb1287.
\end{quote}

\clearpage
\section{Coefficient binding and exact bytes}\label{app:bytes}
The active compiler has
\[
\rho=2^{121},\qquad L=8801,\qquad d=121\cdot8801=1064921,
\qquad B_{\rm native}=2^d.
\]
Its entire finite recipe, transition table, tile table, and controller
proof are included in \file{vendor/}. The unsigned integers are
encoded little-endian and gzip-compressed in \file{data/}.
The raw integer hashes are:
\begin{description}[style=nextline,font=\normalfont\bfseries]
\item[\(\beta_{\rm native}\): \file{beta.bin.gz}, 1,064,921 bits]
\small\ttfamily
6a029139964239dc65d716ccc186b40d9c5923ada0a4116cbb418d84534096ea
\item[\(K_0\): \file{K.bin.gz}, 1,070,004 bits]
\small\ttfamily
6984850dcfc65498b42b633855ea81b8c2cf6d0a84e49005cca10f25557ae395
\item[\(M_C\): \file{MC.bin.gz}, 1,064,921 bits]
\small\ttfamily
7c2b34179cb98c9bea1954f816f93fdc1f60213c68d8e98f1541edb41f0e3a4b
\item[\(M_F\): \file{MF_source.bin.gz}, 1,064,922 bits]
\small\ttfamily
28c55970993394dd7ee32d7299a215fd26b4d70e437a643b60594eee8c7cd3a7
\end{description}
The active source constant is \(\beta_S=2^{2134940}\).
It is distinct from \(\beta_{\rm native}=B_{\rm native}-1\).
The final equation's fixed leaves are specified completely by
these four integers, that explicit power of two, the finite
coefficient recipes, and literal small integers.

\section{A proof dependency map}\label{app:dependencies}
\begin{center}\small
\begin{tabular}{@{}>{\raggedright\arraybackslash}p{42mm}>{\raggedright\arraybackslash}p{49mm}>{\raggedright\arraybackslash}p{58mm}@{}}
\toprule
Step & Statement used & Status in this package\\
\midrule
K215 endpoint &
Blank halting iff a finite ZFC refutation &
Inherited endpoint source and precise theorem name\\[2mm]
Native compiler and positive arithmetic &
The \(31\)-positive-variable system has a solution iff K215 halts &
Inherited ordinary proofs, literal inputs, regenerated coefficient tuple\\[2mm]
Direct residual source &
Common quadratic zero iff the binary/binomial source interface &
New proof, full symbolic coefficient checks, exact data\\[2mm]
Lucas and interpolation bridge &
Source interface iff the specified square/divisibility certificate &
Sun 2026 theorems, same formulas as package 31\\[2mm]
Final mergers and norm &
Certificate iff one equation in ten integers &
Inherited package-31 arguments reproduced in this report\\[2mm]
Degree and circuit transcription &
Exact degree \(237520456\), \(809\) operations &
Independent arithmetic audit and leading-coefficient certificate\\
\bottomrule
\end{tabular}
\end{center}

\clearpage
\begin{thebibliography}{9}
\raggedright
\bibitem{repo}
V.~Reshetnikov and contributors,
\emph{ProveIt: Con(ZFC) Diophantine forms collection},
audited at commit \pin, 23 October 2026.
\url{https://github.com/VladimirReshetnikov/ProveIt/tree/22e75dd6b87b1b199bb08c778d7cf075ae5aae08/SetTheory/Cardinals/docs/reports/hilbert-tenth-problem/con-zfc-diophantine-forms}.
\bibitem{repo31}
ProveIt, package 31, ten-integer source, coefficient compiler, and
structural-weight quintic finalizer. Pinned source copies and notices
are included in this archive. See the collection in \cite{repo}.
\bibitem{sun2021}
Zhi-Wei Sun, \emph{Further results on Hilbert's Tenth Problem},
Science China Mathematics \textbf{64} (2021), 281--306.
Version consulted: arXiv:1704.03504v7.
\url{https://arxiv.org/html/1704.03504v7}.
\bibitem{sun2026}
Zhi-Wei Sun, \emph{Ten unknowns for Hilbert's tenth problem over
the integers}, arXiv:2609.01594v2, 2026 preprint.
Theorems 3.1 and 4.1, Propositions 6.1 and 6.2, and Section 7.
\url{https://arxiv.org/html/2609.01594v2}.
\bibitem{bose}
R.~C.~Bose and S.~Chowla, \emph{Theorems in the additive theory
of numbers}, Commentarii Mathematici Helvetici \textbf{37}
(1962/63), 141--147.
\url{https://doi.org/10.1007/BF02566968}.
\bibitem{nathanson}
M.~B.~Nathanson, \emph{The Bose--Chowla argument for Sidon sets},
Journal of Number Theory \textbf{238} (2022), 133--146.
\url{https://arxiv.org/abs/2104.12711}.
\end{thebibliography}
\end{document}
